\section*{References}
\addcontentsline{toc}{section}{References}
\begingroup
\renewcommand{\section}[2]{}
\makeatletter\renewcommand{\@biblabel}[1]{#1.}\makeatother
\begin{thebibliography}{99}
\setlength{\itemsep}{2pt}

\bibitem{gbd2023ckd}
GBD 2023 CKD Collaborators. \href{https://doi.org/10.1016/S0140-6736(25)01853-7}{Global, regional, and national burden of chronic kidney disease in adults, 1990--2023}. \textit{Lancet}, 406, 2461--2482, 2025.

\bibitem{polat2017}
Polat, H. et al. \href{https://doi.org/10.1007/s10916-017-0703-x}{Diagnosis of chronic kidney disease based on support vector machine by feature selection methods}. \textit{J. Med. Syst.}, 41, 2017.

\bibitem{qin2020}
Qin, J. et al. \href{https://doi.org/10.1109/ACCESS.2019.2963053}{A machine learning methodology for diagnosing chronic kidney disease}. \textit{IEEE Access}, 8, 20991--21002, 2020.

\bibitem{islam2022}
Islam, M.N. et al. \href{https://doi.org/10.1038/s41598-022-15634-4}{Vision transformer and explainable transfer learning models for auto detection of kidney cyst, stone and tumor from CT-radiography}. \textit{Sci. Rep.}, 12, 11440, 2022.

% \bibitem{uci_ckd}
% Rubini, L. et al. \href{https://doi.org/10.24432/C5G020}{Chronic Kidney Disease} [Dataset]. UCI Machine Learning Repository, 2015.
% (cite in Section 3)

\end{thebibliography}
\endgroup
